\addtocontents{toc}{\protect\setcounter{tocdepth}{1}}
\chapter[Factibilidad]{Factibilidad}

Este capítulo evalúa si el proyecto puede llevarse a cabo desde tres
perspectivas: técnica, operativa y económica. Como se trata de un prototipo
académico de salud digital sin fines comerciales, la evaluación económica
valoriza los ahorros de tiempo clínico y de material en lugar de ingresos por
ventas.

\section[Técnica]{Factibilidad técnica}

\subsection{Ambiente de desarrollo}

El desarrollo se realiza en los computadores personales de las estudiantes, con
Flutter SDK y Android Studio, Node.js y PostgreSQL para el backend, Docker
Compose para ejecutar el sistema completo, Git y GitHub para el control de
versiones y Figma para el diseño de la interfaz. Todas estas herramientas son
gratuitas o de código abierto.

\subsection{Ambiente de operación}

La aplicación se instala en teléfonos Android e iOS de gama media. No requiere
equipos adicionales en el hospital ni en el hogar del cuidador.

\subsection{Infraestructura de servidores y datos}

El sistema se compone de una API REST en Node.js y una base de datos PostgreSQL,
ambas de código abierto y empaquetadas en contenedores. Para el despliegue se
cuenta con el servicio Podman Tesistas del servidor que la Universidad del
Bío-Bío dispone para proyectos de título, que cubre la ejecución de los
contenedores y el almacenamiento de los archivos educativos para el volumen de
un prototipo, sin costo para el proyecto. La activación de la cuenta en ese
servidor está pendiente; mientras tanto, el sistema completo se ejecuta
localmente con Docker Compose y se ha probado con la aplicación instalada en un
teléfono Android. Los correos de verificación y recuperación se envían por SMTP
desde una cuenta de Gmail del proyecto.

\begin{table}[H]
\centering
\small
\tablazebra
\caption{Recursos técnicos del proyecto.}
\label{tab:recursos-tecnicos}
\begin{tabularx}{\textwidth}{P{0.95}P{1.60}P{0.45}}
\toprule
\filaencabezado \encab{Recurso} & \encab{Detalle} & \encab{Costo} \\
\midrule
Computadores de desarrollo & Equipos personales de las estudiantes. & \$0 \\
Teléfonos de prueba & Dispositivos Android e iOS de las estudiantes. & \$0 \\
Herramientas de software & Flutter, Node.js, PostgreSQL, Docker, GitHub y Figma
en sus versiones gratuitas o de código abierto. & \$0 \\
Servidor de la universidad (Podman Tesistas) & Ejecución de la API y la base de
datos en contenedores. & \$0 \\
Servicio de correo (Gmail, cuenta del proyecto) & Envío de los correos de
verificación y recuperación. & \$0 \\
\midrule
\filatotal \textbf{Total} & & \textbf{\$0} \\
\bottomrule
\end{tabularx}
\end{table}

\subsection{Capacidades del equipo}

Las estudiantes cuentan con formación en programación, bases de datos,
ingeniería de requerimientos y diseño de interfaces. Las brechas en desarrollo
móvil con Flutter, desarrollo de la API con Node.js y PostgreSQL, y cifrado en el
cliente se cubren con
documentación oficial, apoyo del docente guía y aprendizaje progresivo en cada
sprint. El conocimiento clínico no recae en el equipo de desarrollo, sino en la
contraparte de Enfermería, que define los contenidos y los criterios de alerta.

\subsection{Restricciones técnicas}

El contexto de salud impone dos restricciones. La primera es proteger los datos
sensibles, lo que se resuelve con cifrado AES-256-GCM en el dispositivo y un
servidor que rechaza guardar esos datos en claro. La segunda es no entregar
respuestas fuera del alcance validado, por lo que la orientación se basa en
reglas y respuestas predefinidas, sin inteligencia artificial generativa. Esta
decisión también elimina costos por uso de servicios externos y permite que la
orientación funcione sin conexión.

\section[Operativa]{Factibilidad operativa}

La aplicación está dirigida a cuidadores con distintos niveles de familiaridad
tecnológica y sin formación clínica. Por ello, cada pantalla se centra en una
tarea concreta y usa un lenguaje cotidiano. La implantación no modifica los
procesos del hospital, porque el cuidador sigue asistiendo a sus controles y el
equipo clínico conserva su forma de trabajo. La aplicación solo agrega un
historial ordenado que el cuidador puede mostrar en cada control.

Los impactos esperados se resumen a continuación.

\begin{itemize}
    \item El cuidador accede a información validada en el momento de la duda y
    lleva un registro ordenado del seguimiento.
    \item El paciente se beneficia de una observación más sistemática que
    favorece la detección oportuna de signos de alerta.
    \item El equipo clínico dispone de antecedentes domiciliarios estructurados
    para revisar en cada control.
    \item Como impacto negativo posible, un cuidador podría confiar en exceso en
    la clasificación de la aplicación. Para mitigarlo, cada alerta recuerda que
    la orientación es educativa y recomienda contactar al equipo de salud ante
    cualquier riesgo.
\end{itemize}

El proyecto cuenta con el apoyo de la contraparte de Enfermería de la
Universidad del Bío-Bío, que entrega los contenidos, revisa los criterios de
alerta y valida los textos de la aplicación. Por lo anterior, el proyecto se
considera operativamente factible.

\section[Económica]{Factibilidad económica}

\subsection{Costos}

El costo principal es el esfuerzo de desarrollo. Se consideran dos
desarrolladoras con una dedicación de 15~horas semanales durante 16~semanas y
un valor hora referencial de \$5.231.

\begin{table}[H]
\centering
\small
\tablazebra
\caption{Estimación del costo de desarrollo.}
\label{tab:estimacion-costos}
\begin{tabularx}{\textwidth}{P{1.40}P{0.60}}
\toprule
\filaencabezado \encab{Concepto} & \encab{Valor} \\
\midrule
Valor hora por desarrolladora & \$5.231 \\
Horas semanales por desarrolladora & 15 \\
Semanas de desarrollo & 16 \\
Horas totales por desarrolladora & 240 \\
Desarrolladoras & 2 \\
Horas totales de desarrollo & 480 \\
\midrule
\filatotal \textbf{Costo total de desarrollo} & \textbf{\$2.510.880} \\
\bottomrule
\end{tabularx}
\end{table}

Los computadores, teléfonos, conexión a internet y herramientas ya están
disponibles, por lo que no suman costo. El servidor de la universidad y la
cuenta de correo del proyecto no tienen costo durante la etapa de prototipo.

\begin{table}[H]
\centering
\small
\tablazebra
\caption{Inversión inicial.}
\label{tab:inversion-inicial}
\begin{tabularx}{\textwidth}{P{1.15}P{1.25}P{0.60}}
\toprule
\filaencabezado \encab{Ítem} & \encab{Detalle} & \encab{Costo} \\
\midrule
Desarrollo del prototipo & 480 horas a \$5.231 la hora & \$2.510.880 \\
Herramientas de desarrollo & Versiones gratuitas & \$0 \\
Computadores y teléfonos & Equipos disponibles & \$0 \\
Servidor de la universidad (Podman Tesistas) & Servicio provisto por la universidad & \$0 \\
Servicio de correo (Gmail, cuenta del proyecto) & Cuenta existente & \$0 \\
\midrule
\filatotal \textbf{Inversión inicial} & & \textbf{\$2.510.880} \\
\bottomrule
\end{tabularx}
\end{table}

La mantención anual considera 10~horas por desarrolladora para ajustes menores y
corrección de errores. La revisión del contenido clínico la realiza la
contraparte sin costo para el proyecto.

\begin{table}[H]
\centering
\small
\tablazebra
\caption{Costo anual de mantención.}
\label{tab:mantencion}
\begin{tabularx}{\textwidth}{P{1.50}P{0.50}}
\toprule
\filaencabezado \encab{Concepto} & \encab{Costo anual} \\
\midrule
Servidor de la universidad (Podman Tesistas) & \$0 \\
Servicio de correo (Gmail, cuenta del proyecto) & \$0 \\
Revisión de contenido clínico & \$0 \\
Ajustes técnicos (2 desarrolladoras, 10 h al año, \$5.231 la hora) & \$104.620 \\
\midrule
\filatotal \textbf{Total mantención anual} & \textbf{\$104.620} \\
\bottomrule
\end{tabularx}
\end{table}

\subsection{Beneficios}

Los beneficios tangibles se estiman para un escenario de 50~cuidadores, con el
valor hora de un profesional de enfermería, de \$9.103, como referencia para el
tiempo clínico ahorrado.

\begin{table}[H]
\centering
\small
\tablazebra
\caption{Beneficios tangibles anuales.}
\label{tab:beneficios}
\begin{tabularx}{\textwidth}{P{1.20}P{1.20}P{0.60}}
\toprule
\filaencabezado \encab{Beneficio} & \encab{Cálculo} & \encab{Monto anual} \\
\midrule
Dudas frecuentes resueltas sin llamar al equipo clínico &
50 cuidadores, 4 consultas al año, 0,5~h cada una &
\$910.300 \\
Revisión más rápida de antecedentes en los controles &
50 cuidadores, 10 controles al año, 10 minutos ahorrados en cada uno &
\$758.583 \\
Menos material educativo impreso &
50 cuidadores, 20 hojas al año, \$25 por hoja &
\$25.000 \\
\midrule
\filatotal \textbf{Beneficios anuales} & & \textbf{\$1.693.883} \\
\bottomrule
\end{tabularx}
\end{table}

Entre los beneficios intangibles se cuentan la menor incertidumbre del cuidador,
la detección más oportuna de signos de alerta y una educación homogénea para
todas las familias.

\subsection{Flujo de caja y VAN}

El flujo de caja considera un horizonte de cinco años. El año~0 registra la
inversión inicial y desde el año~1 se suman los beneficios y se descuenta la
mantención.

\begin{table}[H]
\centering
\footnotesize
\tablazebra
\caption{Flujo de caja neto del proyecto.}
\label{tab:flujo-caja}
\begin{tabularx}{\textwidth}{P{1.72}*{6}{>{\raggedleft\arraybackslash\hsize=0.88\hsize}X}}
\toprule
\filaencabezado \encab{Concepto} & \encab{Año 0} & \encab{Año 1} & \encab{Año 2} & \encab{Año 3} & \encab{Año 4} & \encab{Año 5} \\
\midrule
Beneficios (+) & \$0 & \$1.693.883 & \$1.693.883 & \$1.693.883 & \$1.693.883 & \$1.693.883 \\
Mantención (--) & \$0 & -\$104.620 & -\$104.620 & -\$104.620 & -\$104.620 & -\$104.620 \\
Inversión inicial (--) & -\$2.510.880 & \$0 & \$0 & \$0 & \$0 & \$0 \\
\midrule
\filatotal \textbf{Flujo neto} & \textbf{-\$2.510.880} & \textbf{\$1.589.263} & \textbf{\$1.589.263} & \textbf{\$1.589.263} & \textbf{\$1.589.263} & \textbf{\$1.589.263} \\
\bottomrule
\end{tabularx}
\end{table}

Con una tasa de descuento del 10~\% anual, el valor actual neto es:

\begin{equation}
\begin{split}
VAN &= -2.510.880 + \sum_{t=1}^{5} \frac{1.589.263}{(1+0{,}10)^t} \\
    &= \$3.513.677
\end{split}
\label{eq:van}
\end{equation}

La tasa interna de retorno alcanza cerca de un 56~\% y la inversión se recupera
en aproximadamente un año y siete meses. Como el VAN es positivo, el proyecto
resulta económicamente conveniente desde la perspectiva del ahorro social y
operativo.

\section[Conclusión]{Conclusión de la factibilidad}

El proyecto es factible en las tres dimensiones evaluadas. En lo técnico, las
herramientas son gratuitas, están disponibles y el equipo cuenta con las
capacidades necesarias, apoyado por la contraparte en lo clínico. En lo
operativo, la aplicación responde a una necesidad concreta del cuidador sin
alterar los procesos del hospital. En lo económico, el VAN de \$3.513.677 y la
recuperación de la inversión en menos de dos años respaldan el desarrollo.

\addtocontents{toc}{\protect\setcounter{tocdepth}{2}}
